\section{Introduction}

\subsection{The Clinical Challenge: Communication in Aphasia and Polyhandicap}

Patients with aphasia and polyhandicap represent one of the most challenging populations for affective communication. The term "polyhandicap" refers to the combination of severe intellectual disability and motor impairment resulting from non-progressive brain damage occurring during developmental periods. These patients often present with multifocal brain lesions of varying etiologies, atypical functional connectivity, comorbid epilepsy requiring anti-epileptic medications that alter EEG patterns, fluctuating arousal and vigilance states, spasticity generating muscle artifacts, and an inability to provide self-report or follow task instructions \cite{mikolajewska2014noninvasive_bci, schiff2024bci_comms}.

In this context, caregivers and family members must rely on behavioral observation alone to infer internal states such as pain, discomfort, anxiety, or positive affect. However, behavioral expression is often blunted, ambiguous, or absent, leading to potential under-recognition of suffering and missed opportunities for comfort or positive interaction \cite{pan2018emotion_doc}.

\subsection{EEG as a Window to Affective States}

Electroencephalography (EEG) measures scalp potentials generated by large populations of cortical pyramidal neurons. Because it is non-invasive, bedside-compatible, and offers millisecond temporal resolution, EEG is particularly well suited for fragile populations. Even when overt behavior is absent, EEG can reveal preserved cognitive or affective processing, as demonstrated in studies on disorders of consciousness and locked-in syndrome, where brain-computer interfaces (BCIs) have been used to infer covert responses or residual awareness \cite{galiotta2022eeg_bci_doc, he2022bci_doc, lesenfants2015lockedin}.

EEG-based emotion recognition has become a major topic in affective computing, with numerous studies demonstrating that emotional valence and arousal are reflected in changes in oscillatory activity, including modulations of frontal alpha asymmetry, beta and gamma power, and large-scale connectivity patterns \cite{liu2024graph_eeg_emotion, gkintoni2025systematic, zhang2024challenges}. Comprehensive reviews highlight the diversity of feature extraction and classification methods, from classical machine learning to deep neural networks, and the importance of robust signal processing pipelines to handle noise and inter-subject variability \cite{dhuli2022tutorial, pillalamarri2025multimodal}.

\subsubsection{Frontal Alpha Asymmetry: A Key Biomarker}

Frontal alpha asymmetry (FAA) has emerged as one of the most robust biomarkers of emotional valence. The approach-withdrawal model posits that relative right frontal activity is associated with withdrawal-related emotions (fear, disgust), while relative left frontal activity is associated with approach-related emotions (joy, anger) \cite{coan2006frontal}. FAA is typically computed as:
\begin{equation}
\text{FAA} = \log(\alpha_{\text{right}}) - \log(\alpha_{\text{left}})
\end{equation}
where $\alpha_{\text{right}}$ and $\alpha_{\text{left}}$ represent alpha band power (8-13 Hz) over right and left frontal regions.

\subsubsection{Functional Connectivity in Emotional Processing}

Beyond asymmetry, functional connectivity measures have gained increasing attention. Correlation-based connectivity captures the temporal synchrony between brain regions, with studies showing that emotional states modulate large-scale network organization \cite{liu2024graph_eeg_emotion}. Metrics such as mean absolute correlation and its variability provide summary statistics of network state that have proven discriminative in emotion classification tasks. The standard deviation of absolute correlation, in particular, captures dynamic aspects of connectivity that may reflect transitions between emotional states or variations in arousal.

\subsection{The Bitbrain Air System: A Practical Solution for Field Deployment}

For this proof-of-concept study, we selected the Bitbrain Air wireless EEG system with 8 dry electrodes, specifically designed for ease of use in real-world environments. The system features 8 dry, non-gel active electrodes based on silver/silver chloride (Ag/AgCl) technology, arranged according to a reduced 10-20 international system montage covering frontal (F3, F4, Fz), central (C3, C4), temporal (T7, T8), and parietal (Pz) regions \cite{bitbrain2024}. The sampling rate is 256 Hz, providing sufficient temporal resolution to capture relevant EEG frequency bands up to low gamma (40 Hz) while maintaining manageable data volumes for real-time processing.

The system offers several key advantages for field deployment:
\begin{itemize}
    \item \textbf{No gel}: Eliminates lengthy preparation and cleanup (2-3 minutes setup)
    \item \textbf{Comfort}: Lightweight design (90g) reduces skin irritation and discomfort
    \item \textbf{Signal quality}: Active electrodes provide SNR comparable to wet systems up to 40 Hz
    \item \textbf{Mobility}: Wireless design allows recording during natural activities
    \item \textbf{Real-time access}: Python SDK provides low-latency data streaming (< 50 ms)
\end{itemize}

While 8 channels provide less spatial information than high-density systems, they represent an optimal trade-off for field deployment, providing sufficient coverage for frontal asymmetry, central motor/arousal regions, temporal social-emotional areas, and posterior attentional regions.

\subsection{Self-Organizing Maps: A Novel Approach for Patient-Specific Modeling}

Self-Organizing Maps (SOMs) offer an attractive framework for modeling idiosyncratic EEG data. SOMs are unsupervised neural algorithms that map high-dimensional feature vectors onto a low-dimensional grid while preserving topological relationships \cite{kaski1993som_eeg}. They have been successfully applied to EEG for:

\begin{itemize}
    \item Monitoring topographic patterns in clinical EEG \cite{joutsiniemi1995som_eeg}
    \item Seizure detection and classification \cite{elo1992som_epileptic}
    \item Similarity analysis of EEG recordings \cite{jahan2014som_similarity}
    \item BCI classifier structures \cite{som_bci_classifier_2015}
    \item Facial gesture recognition from EEG \cite{kawaguchi2024som_facial}
\end{itemize}

The key advantages of SOMs for our application are:
\begin{enumerate}
    \item \textbf{Unsupervised learning}: They discover structure without requiring labels, which is crucial when annotations are sparse
    \item \textbf{Topology preservation}: Similar feature vectors map to nearby neurons, creating a meaningful organization of brain states
    \item \textbf{Visualization}: The 2D grid provides an intuitive visualization of the patient's "brain state space"
    \item \textbf{Patient-specific}: Each patient develops their own map, respecting individual neural organization
    \item \textbf{Longitudinal stability}: Maps can be updated incrementally as new data becomes available
\end{enumerate}

\subsection{Originality and Contribution to the Literature}

This work presents several original contributions that distinguish it from the existing literature:

\begin{enumerate}
    \item \textbf{First application of SOMs to patient-specific emotion recognition}: While SOMs have been used in EEG analysis for seizure detection and topographic pattern recognition, their application to emotion decoding in a patient-specific framework is novel.
    
    \item \textbf{Integration of dynamic connectivity features}: Most studies focus on static connectivity averages. Our identification of connectivity variability (std\_abs\_corr) as the most discriminative feature (0.138) highlights the importance of dynamic measures, opening new research directions.
    
    \item \textbf{Practical low-channel-count system}: Unlike high-density laboratory systems (32-256 channels), our 8-channel approach demonstrates that clinically useful information can be extracted from a practical, easy-to-deploy system.
    
    \item \textbf{Rapid calibration}: With only 4.7 minutes of recording and 9 annotations, we achieved 77.8\% accuracy, suggesting that patient-specific models can be calibrated much faster than previously thought.
    
    \item \textbf{Interpretable AI}: Unlike deep learning "black boxes," our approach identifies specific, neurophysiologically meaningful features (connectivity variability, theta power, FAA), providing insights into the neural correlates of emotional states.
    
    \item \textbf{Ecological annotation interface}: The Streamlit-based interface allows real-time annotation by non-experts, addressing a critical gap between laboratory research and field deployment.
\end{enumerate}

\subsection{Technological Bottlenecks Addressed}

This proof-of-concept has addressed several key technological challenges:

\begin{itemize}
    \item \textbf{Real-time synchronization}: Latency < 50 ms between EEG acquisition and annotation
    \item \textbf{SOM robustness}: Stable convergence even with limited data (279 windows, 4.7 minutes)
    \item \textbf{Discriminative feature identification}: std\_abs\_corr (0.138) and theta power (0.125) as robust markers
    \item \textbf{Significant classification}: 77.8\% accuracy with only 9 annotations (chance = 16.7\%, p < 0.001)
\end{itemize}

\subsection{Objectives of the Present Study}

This study aims to:
\begin{enumerate}
    \item Validate a complete acquisition-to-classification pipeline on a healthy volunteer
    \item Demonstrate SOM's ability to structure EEG feature space in the context of emotional states
    \item Identify the most discriminative EEG characteristics for emotion recognition
    \item Establish a solid scientific and technological foundation for future field deployment
\end{enumerate}